% Part: normal-modal-logic
% Chapter: tableaux
% Section: countermodels

\documentclass[../../../include/open-logic-section]{subfiles}

\begin{document}

\olfileid{nml}{tab}{cou}
      
\olsection{\usetoken{P}{tableau}માંથી પ્રતિવિધાનો}

\begin{explain}
  પૂર્ણતા પ્રમેયનું પ્રમાણ $\Entails !A$ હોય તો $\Proves !A$
  એટલું જ નથી બતાવતું; $\Entails/ !A$ હોય ત્યારે~$!A$નું
  પ્રતિવિધાન બનાવવાની રીત પણ આપે છે.
  \sourcecorrection{OLMD-071}{મૂળમાં આ પ્રતિવિધાનની શરત
  $\Entails/ A$ લખાઈ છે; બાકીના પરિચ્છેદ અને પૂર્ણતા-દાવા
  મુજબ $!A$નું અનુસરણ ન થવું જોઈએ, તેથી $!A$ રાખ્યું છે.}
  $\Log{K}$ના કિસ્સામાં આ રીત \emph{નિર્ણય-પ્રક્રિયા} છે.
  ધારો કે $\Entails/ !A$. ત્યારે
  \olref[cpl]{prop:complete-tableau}નું પ્રમાણ પૂર્ણ
  !!{tableau} બનાવવાની રીત આપે છે. આ એક-સૂત્રના કિસ્સામાં
  રીત હંમેશા સમાપ્ત થાય છે. $\Log{K}$ના વિધાનસંબંધી નિયમો
  ઓછી જટિલતાવાળાં ઉપસર્ગવાળાં !!{formula}s જ ઉમેરે છે;
  એટલે દરેક ચિહ્નિત સૂત્ર $\sFmla{S}{!A}[\sigma]$ પર
  શાખામાં કોઈ વિધાનસંબંધી નિયમ એક વાર લાગુ કરવો પૂરતો છે.
  નવા ઉપસર્ગો માત્ર
  \iftag{prvBox}{$\TRule{\False}{\Box}$}{}\iftag{notprvBox,notprvDiamond}{}{
    અને }\iftag{prvDiamond}{$\TRule{\True}{\Diamond}$}{}
  \iftag{notprvBox,notprvDiamond}{નિયમ}{નિયમો}થી બને છે; તેમને
  પણ એક જ વાર લાગુ કરવાના હોય છે (અને એક જ નવો ઉપસર્ગ મળે છે).
  \iftag{prvBox}{$\TRule{\True}{\Box}$}{}\iftag{notprvBox,notprvDiamond}{}{
    અને }\iftag{prvDiamond}{$\TRule{\False}{\Diamond}$}{}
  \iftag{notprvBox,notprvDiamond}{નિયમ}{નિયમો} કદાચ અનેક વાર
  લાગુ કરવા પડે, પરંતુ દરેક ઉપસર્ગ માટે એક જ વાર; અને નવા
  ઉપસર્ગો સાંત સંખ્યામાં જ બને છે. તેથી સાંત તબક્કાઓ પછી
  રચનામાં બંધ અથવા પૂર્ણ શાખા મળે છે.

  ખુલ્લી પૂર્ણ શાખાવાળું ટેબ્લો બની જાય પછી
  \olref[cpl]{thm:tableau-completeness}નું પ્રમાણ મૂળના
  ઉપસર્ગવાળાં !!{formula}sના ગણને સંતોષતું સ્પષ્ટ નિદર્શન
  આપે છે. આમ $\Gamma \Entails !A$ હોય તો બંધ !!{tableau}
  મળે છે અને $\Gamma \Proves !A$ થાય છે એટલું જ નહિ;
  બંધ !!{tableau}ને યોગ્ય રીતે શોધતાં જો ``પૂર્ણ''
  !!{tableau} મળે, તો $\Gamma \Entails/ !A$ છે તે જાણીએ
  છીએ અને પ્રતિવિધાન પણ બનાવી શકીએ છીએ.
\end{explain}

\iftag{prvBox}{
\begin{ex}
  આપણે જાણીએ છીએ કે
  $\Proves/ \Box(p \lor q) \lif (\Box p \lor \Box q)$.
  ટેબ્લોનું નિર્માણ આ રીતે શરૂ થાય છે:
  \begin{oltableau}
    [\pFmla{\False}{\Box(p \lor q) \lif (\Box p \lor \Box q)}{1},
      just = \TAss, checked
      [\pFmla{\True}{\Box(p \lor q)}{1},
        just = {\TRule{\False}{\lif}[1]}, 
        [\pFmla{\False}{\Box p \lor \Box q}{1},
          just = {\TRule{\False}{\lif}[1]}, checked
          [\pFmla{\False}{\Box p}{1},
            just = {\TRule{\False}{\lor}[3]}, checked
            [\pFmla{\False}{\Box q}{1},
              just = {\TRule{\False}{\lor}[3]}, checked
              [\pFmla{\False}{p}{1.1}, 
                just = {\TRule{\False}{\Box}[4]}, checked
                [\pFmla{\False}{q}{1.2}, 
                  just = {\TRule{\False}{\Box}[5]}, checked
                ]
              ]
            ]
          ]
        ]
      ]
    ]
  \end{oltableau}
  આ !!{tableau} હજી પૂરું થયું નથી. આગળ ટીકચિહ્ન વિનાની
  એકમાત્ર પંક્તિ જોઈએ: પંક્તિ~$2$ પરનું ઉપસર્ગવાળું
  !!{formula} $\sFmla{\True}{\Box(p \lor q)}[1]$.
  પદ્ધતિસર રચનામાં $\TRule{\True}{\Box}$ નિયમ શાખા પર
  વપરાયેલા દરેક ઉપસર્ગ, એટલે $1.1$ અને $1.2$ બંને માટે
  લાગુ કરવો પડે છે:
  \begin{oltableau}
    [\pFmla{\False}{\Box(p \lor q) \lif (\Box p \lor \Box q)}{1},
      just = \TAss, checked
      [\pFmla{\True}{\Box(p \lor q)}{1},
        just = {\TRule{\False}{\lif}[1]}, 
        [\pFmla{\False}{\Box p \lor \Box q}{1},
          just = {\TRule{\False}{\lif}[1]}, checked
          [\pFmla{\False}{\Box p}{1},
            just = {\TRule{\False}{\lor}[3]}, checked
            [\pFmla{\False}{\Box q}{1},
              just = {\TRule{\False}{\lor}[3]}, checked
              [\pFmla{\False}{p}{1.1}, 
                just = {\TRule{\False}{\Box}[4]}, checked
                [\pFmla{\False}{q}{1.2}, 
                  just = {\TRule{\False}{\Box}[5]}, checked
                  [\pFmla{\True}{p \lor q}{1.1}, 
                    just = {\TRule{\True}{\Box}[2]}
                    [\pFmla{\True}{p \lor q}{1.2}, 
                      just = {\TRule{\True}{\Box}[2]}
                    ]
                  ]
                ]
              ]
            ]
          ]
        ]
      ]
    ]
  \end{oltableau}
  હવે પંક્તિઓ 2, 8 અને 9 પર ટીકચિહ્ન નથી. પરંતુ નવો
  ઉપસર્ગ ઉમેરાયો નથી; તેથી મળતી બધી શાખાઓ પર, તે બંધ ન
  થાય ત્યાં સુધી, પંક્તિઓ~8 અને~9ને
  $\TRule{\True}{\lor}$ લાગુ કરીએ છીએ:
  \begin{oltableau}
    [\pFmla{\False}{\Box(p \lor q) \lif (\Box p \lor \Box q)}{1},
      just = \TAss, checked
      [\pFmla{\True}{\Box(p \lor q)}{1},
        just = {\TRule{\False}{\lif}[1]}, checked
        [\pFmla{\False}{\Box p \lor \Box q}{1},
          just = {\TRule{\False}{\lif}[1]}, checked
          [\pFmla{\False}{\Box p}{1},
            just = {\TRule{\False}{\lor}[3]}, checked
            [\pFmla{\False}{\Box q}{1},
              just = {\TRule{\False}{\lor}[3]}, checked
              [\pFmla{\False}{p}{1.1}, 
                just = {\TRule{\False}{\Box}[4]}, checked
                [\pFmla{\False}{q}{1.2}, 
                  just = {\TRule{\False}{\Box}[5]}, checked
                  [\pFmla{\True}{p \lor q}{1.1}, 
                    just = {\TRule{\True}{\Box}[2]}, checked
                    [\pFmla{\True}{p \lor q}{1.2}, 
                      just = {\TRule{\True}{\Box}[2]}, checked
                      [\pFmla{\True}{p}{1.1},
                        just = {\TRule{\True}{\lor}[8]}, checked, close
                      ]
                      [\pFmla{\True}{q}{1.1},
                        just = {\TRule{\True}{\lor}[8]}, checked
                        [\pFmla{\True}{p}{1.2},
                          just = {\TRule{\True}{\lor}[9]}, checked]
                        [\pFmla{\True}{q}{1.2},
                          just = {\TRule{\True}{\lor}[9]}, checked, close]
                      ]
                    ]
                  ]
                ]
              ]
            ]
          ]
        ]
      ]
    ]
  \end{oltableau}
  એક ખુલ્લી શાખા બાકી રહે છે અને તે પૂર્ણ છે. તેમાંથી
  જગતો $W = \{1, 1.1, 1.2\}$ (ખુલ્લી શાખા પરના ઉપસર્ગો),
  સુલભતા-સંબંધ $R = \{\tuple{1, 1.1}, \tuple{1, 1.2}\}$,
  અને નિમણૂક $V(p) = \{1.2\}$ (કારણ કે પંક્તિ~11 પર
  $\sFmla{\True}{p}[1.2]$ છે) તથા $V(q) = \{1.1\}$
  (કારણ કે પંક્તિ~10 પર $\sFmla{\True}{q}[1.1]$ છે)
  ધરાવતું નિદર્શન વ્યાખ્યાયિત કરીએ છીએ. તે
  \olref{fig:counter-Box}માં દર્શાવ્યું છે; તે
  $\Box(p \lor q) \lif (\Box p \lor \Box q)$નું પ્રતિવિધાન
  છે તે ચકાસી શકો છો.
  \begin{figure}
  \begin{center}
    \begin{tikzpicture}[modal]
      \node[world] (w1) [label={[align=right]right:\mFalse{p}\\ \mFalse{q}}]{$1$}; 
      \node[world] (w2) [label={[align=right]right:\mFalse{p}\\ \mTrue{q}},
        above left=of w1]{$1.1$}; 
      \node[world] (w3) [label={[align=right]right:\mTrue{p}\\ \mFalse{q}},
        above right=of w1] {$1.2$};
      \draw[->] (w1) to (w2);
      \draw[->] (w1) to (w3);
    \end{tikzpicture}
  \end{center}
  \caption{$\Box(p \lor q) \lif (\Box p \lor \Box q)$નું પ્રતિવિધાન.}
  \ollabel{fig:counter-Box}
\end{figure}
\end{ex}          
}
{
\begin{ex}
  આપણે જાણીએ છીએ કે
  $\Proves/ (\Diamond p \land \Diamond q) \lif \Diamond(p
  \land q)$. ટેબ્લોનું નિર્માણ આ રીતે શરૂ થાય છે:
  \begin{oltableau}
    [\pFmla{\False}{(\Diamond p \land \Diamond q) \lif \Diamond(p \land q)}{1},
      just = \TAss, checked
      [\pFmla{\True}{\Diamond p \land \Diamond q}{1},
        just = {\TRule{\False}{\lif}[1]}, checked
        [\pFmla{\False}{\Diamond(p \land q)}{1},
          just = {\TRule{\False}{\lif}[1]}
          [\pFmla{\True}{\Diamond p}{1},
            just = {\TRule{\True}{\land}[2]}, checked
            [\pFmla{\True}{\Diamond q}{1},
              just = {\TRule{\True}{\land}[2]}, checked
              [\pFmla{\True}{p}{1.1}, 
                just = {\TRule{\True}{\Diamond}[4]}, checked
                [\pFmla{\True}{q}{1.2}, 
                  just = {\TRule{\True}{\Diamond}[5]}, checked
                ]
              ]
            ]
          ]
        ]
      ]
    ]
  \end{oltableau}
  આ !!{tableau} હજી પૂરું થયું નથી. આગળ ટીકચિહ્ન વિનાની
  એકમાત્ર પંક્તિ જોઈએ: પંક્તિ~$3$ પરનું ઉપસર્ગવાળું
  !!{formula} $\sFmla{\False}{\Diamond(p \land q)}[1]$.
  પદ્ધતિસર રચનામાં $\TRule{\False}{\Diamond}$ નિયમ
  શાખા પર વપરાયેલા દરેક ઉપસર્ગ, એટલે $1.1$ અને $1.2$
  બંને માટે લાગુ કરવો પડે છે:
  \sourcecorrection{OLMD-072}{મૂળના આ વર્ણનમાં પંક્તિ~3ને
  સત્ય ડાયમંડ અને નિયમને સત્ય ડાયમંડનો નિયમ કહ્યા છે;
  પ્રથમ તથા અનુગામી વૃક્ષોમાં પંક્તિ~3 અસત્ય ડાયમંડ છે,
  જેના માટે વપરાયેલા બંને ઉપસર્ગે અસત્ય-ડાયમંડ નિયમ લાગે છે.}
  \sourcecorrection{OLMD-073}{મૂળના બીજા વૃક્ષનું મૂળ
  પ્રથમ અને ત્રીજા વૃક્ષથી જુદા ક્રમનો અનુબંધ દર્શાવે છે,
  જ્યારે તેની બધી આગળની પંક્તિઓ પ્રથમ વૃક્ષના અનુબંધને
  અનુસરે છે. બીજા વૃક્ષનું મૂળ બંને સાથે સુસંગત કર્યું છે.}
  \begin{oltableau}
    [\pFmla{\False}{(\Diamond p \land \Diamond q) \lif \Diamond(p \land q)}{1},
      just = \TAss, checked
      [\pFmla{\True}{\Diamond p \land \Diamond q}{1},
        just = {\TRule{\False}{\lif}[1]}, checked
        [\pFmla{\False}{\Diamond (p \land q)}{1},
          just = {\TRule{\False}{\lif}[1]}
          [\pFmla{\True}{\Diamond p}{1},
            just = {\TRule{\True}{\land}[2]}, checked
            [\pFmla{\True}{\Diamond q}{1},
              just = {\TRule{\True}{\land}[2]}, checked
              [\pFmla{\True}{p}{1.1}, 
                just = {\TRule{\True}{\Diamond}[4]}, checked
                [\pFmla{\True}{q}{1.2}, 
                  just = {\TRule{\True}{\Diamond}[5]}, checked
                  [\pFmla{\False}{p \land q}{1.1}, 
                    just = {\TRule{\False}{\Diamond}[3]}
                    [\pFmla{\False}{p \land q}{1.2}, 
                      just = {\TRule{\False}{\Diamond}[3]}
                    ]
                  ]
                ]
              ]
            ]
          ]
        ]
      ]
    ]
  \end{oltableau}
  હવે પંક્તિઓ 3, 8 અને 9 પર ટીકચિહ્ન નથી. પરંતુ નવો
  ઉપસર્ગ ઉમેરાયો નથી; તેથી મળતી બધી શાખાઓ પર, તે બંધ ન
  થાય ત્યાં સુધી, પંક્તિઓ~8 અને~9ને
  $\TRule{\False}{\land}$ લાગુ કરીએ છીએ:
  \begin{oltableau}
    [\pFmla{\False}{(\Diamond p \land \Diamond q) \lif \Diamond(p \land q)}{1},
      just = \TAss, checked
      [\pFmla{\True}{\Diamond p \land \Diamond q}{1},
        just = {\TRule{\False}{\lif}[1]}, checked
        [\pFmla{\False}{\Diamond(p \land q)}{1},
          just = {\TRule{\False}{\lif}[1]}
          [\pFmla{\True}{\Diamond p}{1},
            just = {\TRule{\True}{\land}[2]}, checked
            [\pFmla{\True}{\Diamond q}{1},
              just = {\TRule{\True}{\land}[2]}, checked
              [\pFmla{\True}{p}{1.1}, 
                just = {\TRule{\True}{\Diamond}[4]}, checked
                [\pFmla{\True}{q}{1.2}, 
                  just = {\TRule{\True}{\Diamond}[5]}, checked
                  [\pFmla{\False}{p \land q}{1.1}, 
                    just = {\TRule{\False}{\Diamond}[3]}, checked
                    [\pFmla{\False}{p \land q}{1.2}, 
                      just = {\TRule{\False}{\Diamond}[3]}, checked
                      [\pFmla{\False}{p}{1.1},
                        just = {\TRule{\False}{\land}[8]}, checked, close
                      ]
                      [\pFmla{\False}{q}{1.1},
                        just = {\TRule{\False}{\land}[8]}, checked
                        [\pFmla{\False}{p}{1.2},
                          just = {\TRule{\False}{\land}[9]}, checked]
                        [\pFmla{\False}{q}{1.2},
                          just = {\TRule{\False}{\land}[9]}, checked, close]
                      ]
                    ]
                  ]
                ]
              ]
            ]
          ]
        ]
      ]
    ]
  \end{oltableau}
  એક ખુલ્લી શાખા બાકી રહે છે અને તે પૂર્ણ છે. તેમાંથી
  જગતો $W = \{1, 1.1, 1.2\}$ (ખુલ્લી શાખા પરના ઉપસર્ગો),
  સુલભતા-સંબંધ $R = \{\tuple{1, 1.1}, \tuple{1, 1.2}\}$,
  અને નિમણૂક $V(p) = \{1.1\}$ (કારણ કે પંક્તિ~6 પર
  $\sFmla{\True}{p}[1.1]$ છે) તથા $V(q) = \{1.2\}$
  (કારણ કે પંક્તિ~7 પર $\sFmla{\True}{q}[1.2]$ છે)
  ધરાવતું નિદર્શન વ્યાખ્યાયિત કરીએ છીએ.
  \sourcecorrection{OLMD-074}{મૂળમાં પંક્તિ~7ના સત્ય
  $q$નું સ્થાન $1.1$ લખાયું છે; વૃક્ષ અને નિમણૂક
  $V(q)=\{1.2\}$ મુજબ તે $1.2$ છે.}
  તે \olref{fig:counter-Diamond}માં દર્શાવ્યું છે;
  $(\Diamond p \land \Diamond q) \lif \Diamond (p \land q)$નું
  પ્રતિવિધાન છે તે ચકાસી શકો છો.
  \begin{figure}
  \begin{center}
    \begin{tikzpicture}[modal]
      \node[world] (w1) [label={[align=right]right:\mFalse{p}\\ \mFalse{q}}]{$1$}; 
      \node[world] (w2) [label={[align=right]right:\mTrue{p}\\ \mFalse{q}},
        above left=of w1]{$1.1$}; 
      \node[world] (w3) [label={[align=right]right:\mFalse{p}\\ \mTrue{q}},
        above right=of w1] {$1.2$};
      \draw[->] (w1) to (w2);
      \draw[->] (w1) to (w3);
    \end{tikzpicture}
  \end{center}
  \caption{$(\Diamond p \land \Diamond q) \lif
    \Diamond (p \land q)$નું પ્રતિવિધાન.}
  \ollabel{fig:counter-Diamond}
\end{figure}
\end{ex}          
}

\end{document}
